% Part: computability
% Chapter: recursive-functions
% Section: halting-problem

\documentclass[../../../include/open-logic-section]{subfiles}

\begin{document}

\olfileid{cmp}{rec}{hlt}
\olsection{અટકવાની સમસ્યા}

સામાન્ય રીતે \emph{અટકવાની સમસ્યા} એટલે સંગણનીય વિધેયનું
વર્ણન~$e$ (દાખલા તરીકે કાર્યક્રમ) અને સંખ્યા~$n$ આપેલાં હોય,
ત્યારે ઇનપુટ~$n$ પર તે વિધેયની ગણતરી અટકે છે કે નહીં, એટલે
પરિણામ આપે છે કે નહીં, તે નક્કી કરવાની સમસ્યા. એલન ટ્યુરિંગે
પ્રસિદ્ધ રીતે સાબિત કર્યું કે આ સમસ્યા પોતે સંગણનીય વિધેયથી
ઉકેલી શકાતી નથી; એટલે નીચેનું વિધેય સંગણનીય નથી:
\[
h(e, n) =
\begin{cases}
1 & \text{જો ગણતરી $e$ ઇનપુટ $n$ પર અટકે}\\
0 & \text{અન્યથા,}
\end{cases}
\]

આંશિક પુનરાવર્તી વિધેયોના સંદર્ભમાં કાર્યક્રમના વર્ણનનું કામ
ક્લીનીના સામાન્ય સ્વરૂપ પ્રમેયમાં મળતો સૂચક~$e$ કરી શકે છે.
$f$ આંશિક પુનરાવર્તી વિધેય હોય, તો સામાન્ય સ્વરૂપ પ્રમેયનું
સમીકરણ સંતોષતો કોઈ પણ $e$ એ~$f$નો સૂચક છે. સંખ્યા~$e$ આપેલી
હોય, તો સામાન્ય સ્વરૂપ પ્રમેય મુજબ
\[
\cfind{e}(x) \simeq U(\mu s \; T(e, x, s))
\]
આંશિક પુનરાવર્તી છે; અને દરેક આંશિક પુનરાવર્તી $f\colon \Nat
\to \Nat$ માટે એવો $e \in \Nat$ છે કે બધા~$x \in \Nat$ માટે
$\cfind{e}(x) \simeq
f(x)$. હકીકતમાં, આવા દરેક $f$ને એક નહીં પણ અનંત એવા~$e$
મળે છે. \emph{અટકવાનું વિધેય}~$h$ આમ વ્યાખ્યાયિત છે:
\[
h(e, x) =
\begin{cases}
1 & \text{જો $\cfind{e}(x) \fdefined$ હોય}\\
0 & \text{અન્યથા.}
\end{cases}
\]
નોંધો કે $\cfind{e}(x) \fundefined$ હોય ત્યારે $h(e, x) = 0$
છે; અને~$e$ બિલકુલ કોઈ આંશિક પુનરાવર્તી વિધેયનો સૂચક ન હોય
ત્યારે પણ એવું જ છે.

\begin{thm}
\ollabel{thm:halting-problem}
અટકવાનું વિધેય~$h$ આંશિક પુનરાવર્તી નથી.
\end{thm}

\begin{proof}
જો $h$ આંશિક પુનરાવર્તી હોત, તો નીચે મુજબ વિધેય વ્યાખ્યાયિત
કરી શકાત:
\[
d(y) =
\begin{cases}
1 & \text{જો $h(y, y) = 0$ હોય}\\
\umin{x}{x \neq x} & \text{અન્યથા.}
\end{cases}
\]
કોઈ સંખ્યા $x$ એ $x \neq x$ સંતોષતી નથી. તેથી
$\umin{x}{x \neq
x}$ અસ્તિત્વમાં નથી અને $h(y,y) \neq 0$ હોય ત્યારે અને માત્ર
ત્યારે જ $d(y) \fundefined$ છે. આ વ્યાખ્યામાંથી મળે છે:
\begin{enumerate}
\item $\cfind{y}(y) \fundefined$ હોય અથવા $y$ એ આંશિક
  પુનરાવર્તી વિધેયનો સૂચક ન હોય ત્યારે અને માત્ર ત્યારે જ
  $d(y) \fdefined$ છે.
\item $\cfind{y}(y) \fdefined$ હોય ત્યારે અને માત્ર ત્યારે જ
  $d(y) \fundefined$ છે.
\end{enumerate}
જો $h$ આંશિક પુનરાવર્તી હોત, તો $d$ પણ આંશિક પુનરાવર્તી હોત.
તેથી ક્લીનીના સામાન્ય સ્વરૂપ પ્રમેય મુજબ તેને સૂચક~$e_d$ મળે.
$h(e_d, e_d)$નું મૂલ્ય વિચારો. તેના બે શક્ય કિસ્સા છે: $0$ અને~$1$.
\begin{enumerate}
\item જો $h(e_d, e_d) = 1$ હોય, તો $\cfind{e_d}(e_d)
  \fdefined$. પરંતુ $\cfind{e_d} \simeq d$ છે; અને
  $h(e_d, e_d) = 0$ હોય ત્યારે અને માત્ર ત્યારે જ $d(e_d)$
  વ્યાખ્યાયિત છે. તેથી $h(e_d, e_d) \neq 1$.
\item જો $h(e_d, e_d) = 0$ હોય, તો કાં તો $e_d$ આંશિક પુનરાવર્તી
  વિધેયનો સૂચક નથી, અથવા છે અને $\cfind{e_d}(e_d)
  \fundefined$. પરંતુ ફરી $\cfind{e_d} \simeq d$ છે; અને
  $\cfind{e_d}(e_d) \fdefined$ હોય ત્યારે અને માત્ર ત્યારે જ
  $d(e_d)$ અવ્યાખ્યાયિત છે.
\end{enumerate}
અંતે એવું નીકળે છે કે $e_d$ આંશિક પુનરાવર્તી વિધેયનો સૂચક હોઈ
જ ન શકે. પરંતુ જો $h$ આંશિક પુનરાવર્તી હોત, તો $d$ પણ હોત;
એટલે $e_d$ને તેનો સૂચક માનતી આપણી વ્યાખ્યા માન્ય હોત. તેથી
$h$ આંશિક પુનરાવર્તી હોઈ શકે નહીં.
\end{proof}

\end{document}
